\documentclass[11pt]{article}
\usepackage{amsmath,amssymb,amsthm}
\usepackage[margin=1in]{geometry}

\newtheorem{theorem}{Theorem}
\newtheorem{remark}{Remark}

\begin{document}

\title{Local Convergence of the Exceptional-Horizon\\
Logarithmic Frobenius Expansion for the $\ell=2$ Axial GfE System}
\date{7 August 2026}
\maketitle

\paragraph{Scope.}
This theorem upgrades the previously certified all-orders formal
exceptional-horizon recurrence to a convergent local logarithmic
Frobenius solution.  It does not establish global analytic
continuation, global black-hole boundary conditions, or physical
instability.

\begin{theorem}[Exceptional-horizon local convergence]
Fix $M>0$ and $\beta\neq0$.  At
\[
  \omega=\frac{i}{4M},
  \qquad
  \alpha=\frac12,
\]
consider the formally certified expansion
\[
  U(x)=x^{1/2}
  \left[
    u_0+\sum_{n\ge1}x^n
    \bigl(u_n+v_n\log x\bigr)
  \right],
  \qquad x=r-2M.
\]
There exist constants $C,Q>0$, depending on the fixed parameters and
the chosen low-order free data, such that
\[
  \|u_n\|+\|v_n\|\le C Q^n
\]
for all $n$.

Consequently there is a radius $\rho>0$ for which
\[
  \sum_{n\ge0}u_nx^n,
  \qquad
  \sum_{n\ge1}v_nx^n
\]
converge absolutely for $|x|<\rho$.  On any fixed local branch of
$\log x$, the resulting expression is therefore a genuine solution of
the axial GfE equations on the punctured neighborhood
\[
  0<|x|<\rho .
\]
For the physical exterior ray this gives a genuine local solution for
$0<x<\rho$.

\paragraph{Proof.}
After the certified horizon row scalings, the exact Euler--Frobenius
symbol has common analytic denominator
\[
  D(x)=\frac{(2M+x)^9}{512M^9},
  \qquad D(0)=1,
\]
and clearing $D$ produces a recurrence of finite lag
\[
  J=10.
\]
The symbol is polynomial of degree at most four in the Frobenius
exponent.  The rank-one compatibility reduction from the all-orders
formal theorem gives a two-component logarithmic state recurrence
whose rational transfer coefficients have no poles at any integer
$n\ge3$.  Every such coefficient has degree at infinity at most zero.
Hence all logarithmic transfer matrices are uniformly bounded:
there is $B<\infty$ such that their norms are bounded by $B$ for every
integer $n\ge3$.

Choose $Q$ sufficiently large that
\[
  \sum_{j=1}^{10} BQ^{-j}<1.
\]
The finite set of initial coefficients can be absorbed into a constant
$C$.  Induction in the finite-lag recurrence then yields
\[
  \|v_n\|\le C Q^n .
\]

The non-logarithmic recurrence has the same bounded homogeneous
transfer.  Its additional source is generated by differentiating the
Euler symbol with respect to the Frobenius exponent.  The certified
derivative symbol has exponent degree at most three, so its finite-lag
forcing is bounded by a polynomial in $n$ times the already geometric
logarithmic sequence.  For every $Q_1>Q$ and every fixed polynomial
$P$ there is $C_1<\infty$ with
\[
  P(n)Q^n\le C_1Q_1^n .
\]
Increasing the geometric base once more therefore closes the same
finite-lag induction for the ordinary coefficients and gives
\[
  \|u_n\|\le C_2 Q_2^n .
\]
Combining the two estimates proves the stated coefficient bound.

The two power series consequently converge absolutely for
$|x|<Q_2^{-1}$.  Since the differential-equation coefficients are
analytic at the horizon after the certified regular-singular row
scalings, termwise differentiation is valid on compact subsets of the
punctured convergence disk.  The formal recurrence therefore implies
the differential equations identically there.
\end{theorem}

\begin{remark}[Certified recurrence structure]
The symbolic verifier obtains
\[
  D(x)=\frac{(2M+x)^9}{512M^9},
  \qquad J=10,
\]
maximum exponent degree four, exponent-derivative degree three, and
maximum logarithmic transfer degree zero.  Thus the convergence proof
does not depend on extending finite-order coefficient calculations.
\end{remark}

\begin{remark}[Boundary]
No optimal value of $\rho$ is asserted.  In particular, this result
does not yet prove continuation throughout the exterior domain,
satisfaction of a condition at infinity, existence of a global mode,
or a physical instability.
\end{remark}

\end{document}
